% Part: set-theory
% Chapter: z
% Section: separation

\documentclass[../../../include/open-logic-section]{subfiles}

\begin{document}

\olfileid{sth}{z}{sep}
\olsection{পৃথকীকরণ}

অনানুষ্ঠানিক ধর্মনির্দেশে সেট-গঠনের বদলে আমরা প্রথমে
নিচের নীতিটি গ্রহণ করি:

\begin{axiom}[পৃথকীকরণ ছক] প্রত্যেক সূত্র~$\phi(x)$-এর জন্য
এটি একটি স্বতঃসিদ্ধ: যেকোনো~$A$-এর ক্ষেত্রে
$\Setabs{x \in A}{\phi(x)}$ সেটটি আছে।
\end{axiom}

লক্ষ করুন, এটি একটি মাত্র স্বতঃসিদ্ধ নয়;
স্বতঃসিদ্ধগুলির একটি \emph{ছক}। পৃথকীকরণের স্বতঃসিদ্ধ
\emph{অসীমসংখ্যক}: প্রত্যেক সূত্র~$\phi(x)$-এর জন্য
একটি। ছকটি সমানভাবেই নিচের রূপে লেখা যায় (এবং সাধারণত
এভাবেই লেখা হয়):

\begin{defish}
``$S$'' নেই এমন যেকোনো সূত্র~$\phi(x)$-এর জন্য
নিচেরটি একটি স্বতঃসিদ্ধ:
\[
  \forall A \exists S \forall x(x \in S \liff (\phi(x) \land x \in A)).
\]
\end{defish}

% BN-SRC-434: অংশের সূচনায় পরামিতির উল্লেখিত রীতিটি বলা নেই।
এই \olref[sth][][]{part}-এ অনুসৃত রীতি অনুযায়ী,
পৃথকীকরণের স্বতঃসিদ্ধগুলিতে
সূত্র~$\phi$-তে পরামিতি থাকতে পারে।\footnote{এর অর্থ কী,
তা জানতে \olref[sfr][infinite][induction]{natinductionschema}-এর
ঠিক পরের আলোচনা দেখুন।}

সেটের সঞ্চয়ী-পর্যায়ক্রমিক ধারণা থেকেই পৃথকীকরণের
যথার্থতা সরাসরি বোঝা যায়। কেন, তা দেখতে $A$-কে
একটি সেট ধরি। তাহলে $A$ কোনো পর্যায়~$S$-এ গঠিত
হয়েছে (\stageshier{} অনুযায়ী)। $A$ ওই পর্যায়~$S$-এ
গঠিত হওয়ায় $A$-এর সব সদস্য $S$-এর আগেই গঠিত হয়েছে
(\stagesacc{} অনুযায়ী)। এবার বিশেষভাবে $A$-এর সেই
সদস্য-সেটগুলি বিবেচনা করি, যেগুলি $\phi$-ও পরিতৃপ্ত করে;
স্পষ্টত সেগুলিও পর্যায়~$S$-এর আগে গঠিত হয়েছে। তাই
পর্যায়~$S$-এ সেগুলিকে নিয়ে
$\Setabs{x \in A}{\phi(x)}$ সেটটিও গঠিত হয়
(\stagesacc{} অনুযায়ী)।

অনানুষ্ঠানিক ধর্মনির্দেশে সেট-গঠনের বিপরীতে, এতে
রাসেলের কূটাভাস এড়ানো যায়। আমরা সরাসরি
$\Setabs{x}{x \notin x}$ সেটটির অস্তিত্ব দাবি করতে
পারি না। বরং কোনো সেট~$A$ \emph{দেওয়া থাকলে}
$R_A = \Setabs{x \in A}{x \notin x}$ সেটটির অস্তিত্ব
দাবি করতে পারি। কিন্তু তাতে শুধু $R_A \notin R_A$
এবং $R_A \notin A$ প্রমাণ হয়; এর কোনোটিই উদ্বেগের
কারণ নয়।

তবে পৃথকীকরণের একটি তাৎক্ষণিক ও লক্ষণীয় ফল আছে:

\begin{thm}\ollabel{thm:NoUniversalSet}
কোনো \emph{সার্বিক} সেট নেই; অর্থাৎ,
$\Setabs{x}{x = x}$ সেটটির অস্তিত্ব নেই।
\end{thm}

\begin{proof}
বিরোধপ্রমাণে ধরি, $V$ একটি সার্বিক সেট। তাহলে
পৃথকীকরণ অনুযায়ী $R =
\Setabs{x \in V}{x \notin x} = \Setabs{x}{x \notin x}$
আছে। এটি রাসেলের কূটাভাসের বিরোধী।
\end{proof}

সার্বিক সেটের অনস্তিত্ব—বরং সেটের ক্রমধারাটি যে
উন্মুক্ত থেকে যায়—সঞ্চয়ী-পর্যায়ক্রমিক ধারণার
সবচেয়ে মৌলিক বিষয়গুলির একটি। তাই স্বাভাবিক বোধে
দেখা ভালো যে, অন্য পথেও এই সিদ্ধান্তে পৌঁছনো যায়।
সার্বিক সেটকে নিজেরই !!a{element} হতে হবে। কিন্তু
আমাদের সঞ্চয়ী-পর্যায়ক্রমিক ধারণায় প্রত্যেক সেট তার
সব !!{element}s যে পর্যায়গুলিতে প্রথম এসেছে, সেগুলির
পরের প্রথম পর্যায়ে নিজে প্রথম দেখা দেয়। এর অর্থ,
\emph{কোনো} সেট নিজের সদস্য নয়। কারণ নিজেকে সদস্য
হিসেবে ধারণকারী সেটটিকে সে প্রথম যে পর্যায়ে এসেছে,
তার অব্যবহিত পরেই আবার প্রথম আসতে হতো—যা অসম্ভব।
(\olref[sth][spine][rank]{defnsetrank}-এ সেটের
``স্তর'' ধারণা দিয়ে এই ব্যাখ্যাকে কীভাবে আরও কঠোর
করা যায়, তা দেখব। তবে তার আগে আরও কয়েকটি
স্বতঃসিদ্ধ স্থির করতে হবে।)

পৃথকীকরণ ও সদস্যভিত্তিক সমতার আরও কয়েকটি ফল এই।

\begin{prop}\ollabel{prop:emptyexists}
যদি কোনো সেট থাকে, তবে $\emptyset$-ও আছে।
\end{prop}

\begin{proof}
$A$ একটি সেট হলে, পৃথকীকরণ অনুযায়ী
$\emptyset = \Setabs{x \in A}{x \neq x}$ আছে।
\end{proof}

\begin{prop}
যেকোনো সেট $A$ ও~$B$-এর জন্য $A \setminus B$ আছে।
\end{prop}

\begin{proof}
পৃথকীকরণ অনুযায়ী
$A \setminus B = \Setabs{x \in A}{x \notin B}$ আছে।
\end{proof}

দেখা যায়, (প্রায়) যেকোনো সংগ্রহের ছেদও আছে:

\begin{prop}\ollabel{prop:intersectionsexist}
যদি $A \neq \emptyset$ হয়, তবে
$\bigcap A = \Setabs{x}{(\forall y \in A)x \in y}$ আছে।
\end{prop}

\begin{proof}
ধরি $A \neq \emptyset$; তাহলে কোনো $c \in A$
আছে। এখন $\bigcap A =
\Setabs{x}{(\forall y \in A)x \in y} = \Setabs{x \in c}{(\forall y \in
A)x \in y}$; পৃথকীকরণ অনুযায়ী এই সেটটি আছে।
\end{proof}

তবে $A \neq \emptyset$ শর্তটি লক্ষ করুন:
$\bigcap \emptyset$ শূন্যার্থে সার্বিক সেট হয়ে যেত,
যা \olref{thm:NoUniversalSet}-এর বিরোধী।

\end{document}
